\chapter{Pierwszy program}

\section{Po co zaczynac od najmniejszego mozliwego przykladu}
Kiedy poznajemy nowy jezyk, bardzo latwo jest od razu zobaczyc duzy przyklad,
przestraszyc sie liczby elementow i stracic intuicje. Dlatego zaczniemy od programu,
ktory robi tylko jedna rzecz: generuje pojedyncza strone glowna.

Oto minimalny przyklad:

\begin{lstlisting}[style=moonchunk,caption={Minimalny program w MoonChunk}]
chunk "Main" {
  output: "./dist";

  page "" {
    content {
      <h1>Witaj w MoonChunk</h1>
      <p>To jest moja pierwsza strona.</p>
    };
  };
};

moon(Main);
\end{lstlisting}

Jest tu zaledwie kilka elementow, ale juz teraz widzimy prawie caly podstawowy model
pracy z jezykiem.

\section{Czytanie programu od gory do dolu}
Przejdzmy przez kod linia po linii.

\subsection*{\mc{chunk "Main" \{ ... \};}}
\mc{chunk} jest podstawowym kontenerem logiki i tresci w MoonChunk. Mozna o nim
myslec jak o module albo opakowaniu na fragment programu. Wewnatrz \mc{chunk}-a
umieszcza sie instrukcje zwiazane z generowaniem stron, definicje funkcji, zmienne,
metadata oraz importowane fragmenty przez \mc{@include}.

Nazwa \mc{"Main"} jest nazwa logiczna. Nie jest to nazwa pliku HTML. To nazwa
obiektu, do ktorego potem odwolujemy sie przez \mc{moon(Main)}.

\subsection*{\mc{output: "./dist";}}
Instrukcja \mc{output} mowi, gdzie maja trafic wygenerowane pliki. W tym przykladzie
wyniki zostana zapisane w katalogu \mc{dist}. Jesli wygenerujemy strone glowna,
prawdopodobnie zobaczymy tam plik \mc{index.html}.

\subsection*{\mc{page "" \{ ... \};}}
\mc{page} oznacza pojedyncza strone, ktora ma zostac wygenerowana. Pusty napis
\mc{""} oznacza trase glowna. W praktyce przeklada sie to na \mc{index.html}.

To bardzo wazny wzorzec:

\begin{itemize}
  \item \mc{page ""} to strona glowna,
  \item \mc{page "about"} to zwykle \mc{about.html},
  \item \mc{page "blog/post-1"} to zwykle \mc{blog/post-1.html}.
\end{itemize}

\subsection*{\mc{content \{ ... \};}}
Blok \mc{content} przechowuje HTML, ktory trafi do glownej czesci dokumentu.
To miejsce, w ktorym najczesciej wpisuje sie znaczniki takie jak \mc{<h1>}, \mc{<p>},
\mc{<section>}, \mc{<article>} czy \mc{<ul>}.

\subsection*{\mc{moon(Main);}}
Ta linia uruchamia wykonanie wybranego \mc{chunk}-a. Bez niej program moze zostac
poprawnie sparsowany, ale nic nie zostanie wygenerowane, poniewaz nie wskazalismy
punktu wejscia.

\section{Dlaczego \mc{moon(...)} jest takie wazne}
W wielu jezykach program zaczyna sie automatycznie od jednej, uprzywilejowanej funkcji.
MoonChunk dziala troche inaczej. Mozesz miec w jednym projekcie wiele \mc{chunk}-ow,
a \mc{moon(...)} jawnie wskazuje, ktory z nich ma zostac wykonany.

Dzieki temu:

\begin{itemize}
  \item latwo pracowac na roznych wariantach projektu,
  \item mozna generowac kilka zestawow stron w jednym przebiegu,
  \item importowane fragmenty nie uruchamiaja sie same z siebie,
  \item kod staje sie bardziej przewidywalny.
\end{itemize}

To niepozorna konstrukcja, ale bardzo porzadkuje duze projekty.

\section{Wersja rozbudowana o zmienne}
Dodajmy teraz dwie proste zmienne i wykorzystajmy je w HTML:

\begin{lstlisting}[style=moonchunk,caption={Pierwszy program z danymi dynamicznymi}]
chunk "Main" {
  output: "./dist";

  page "" {
    let title = "Moja pierwsza strona";
    let intro = "Ten tekst pochodzi ze zmiennej.";

    content {
      <h1>{title}</h1>
      <p>{intro}</p>
    };
  };
};

moon(Main);
\end{lstlisting}

W tym przykladzie \mc{title} i \mc{intro} nie sa wpisane na sztywno w HTML-u.
Najpierw otrzymuja wartosc, a potem sa wstawiane do tresci strony przez zapis
\mc{\{title\}} i \mc{\{intro\}}.

\section{Co dzieje sie podczas wykonania}
Gdy MoonChunk uruchamia taki program, mozna to opisac w nastepujacych krokach:

\begin{enumerate}
  \item odczytuje plik \mc{.mncnk},
  \item rozpoznaje \mc{chunk}-i, strony, zmienne i pozostala logike,
  \item wykonuje \mc{moon(Main)},
  \item wchodzi do wnętrza \mc{chunk "Main"},
  \item dochodzi do strony \mc{page ""},
  \item tworzy zmienne \mc{title} i \mc{intro},
  \item sklada HTML z bloku \mc{content},
  \item zapisuje wynik do pliku wyjsciowego.
\end{enumerate}

Jesli rozumiesz te osiem punktow, rozumiesz juz podstawowy cykl zycia programu
MoonChunk.

\section{Najwazniejsza roznica miedzy kodem a wynikiem}
Dobrze jest rozdzielic w glowie dwie rzeczy:

\begin{itemize}
  \item \important{kod MoonChunk} -- to instrukcje, ktore piszesz,
  \item \important{wygenerowany HTML} -- to finalny wynik po wykonaniu instrukcji.
\end{itemize}

Na przyklad taki kod:

\begin{lstlisting}[style=moonchunk,caption={Kod z interpolacja wartosci}]
let name = "Alicja";
content {
  <p>Czesc, {name}!</p>
};
\end{lstlisting}

moze zamienic sie na wynikowy HTML:

\begin{lstlisting}[style=basecode,caption={Przykladowy wynik po wykonaniu}]
<p>Czesc, Alicja!</p>
\end{lstlisting}

Ta umiejetnosc rozrozniannia etapu \emph{pisania programu} od etapu \emph{wygenerowania
wyniku} jest kluczowa w calej dalszej pracy.

\section{Typowe pomylki w pierwszym programie}
Na poczatku najczesciej trafiaja sie cztery rodzaje bledow:

\begin{itemize}
  \item brak srednika \mc{;},
  \item brak zamykajacego nawiasu klamrowego,
  \item brak \mc{moon(Main);},
  \item wpisanie wartosci dynamicznej bez nawiasow klamrowych w \mc{content}.
\end{itemize}

Przykladowo, jesli napiszesz:

\begin{lstlisting}[style=moonchunk,caption={Zly zapis wartosci dynamicznej}]
content {
  <h1>title</h1>
};
\end{lstlisting}

na stronie zobaczysz doslownie slowo \mc{title}, a nie wartosc zmiennej. Aby wstawic
zawartosc zmiennej, potrzebny jest zapis \mc{\{title\}}.

\section{Co warto zapamietac po tym rozdziale}
Po pierwszym programie powinna zostac Ci w glowie bardzo prosta mapa:

\begin{enumerate}
  \item \mc{chunk} grupuje logike,
  \item \mc{page} opisuje generowana strone,
  \item \mc{content} zawiera HTML,
  \item \mc{moon(...)} uruchamia wybrany punkt wejscia,
  \item \mc{output} wskazuje katalog wynikowy.
\end{enumerate}

Na tych pieciu elementach zbudujemy wszystko, co pojawi sie dalej.
